\section{Pour aller plus loin}

\subsection{Pré-traitements}

\begin{enumerate}
	\item Répétez la procédure pour deux autres classes (chiffres 6 et 8). Est-il plus difficile de séparer ces classes? La courbe d'erreur de validation est-elle différente de la précédente?
	\item Prenez deux photos (noir et blanc, carré) avec votre smartphone, représentant les deux classes (6 et 8), ou utilisez les photos fournies dans les fichiers du tp. Classifiez-les après les étapes de pré-traitement suivantes:
	\begin{itemize}[label=-]
		\item Recadrage
		\item Conversion en négatif (blanc sur fond noir)
		\item Correction du contraste -- il faut normaliser les données de telle façon que les valeurs minimales et maximales deviennent 0 et 255: $f_{i,j} = \frac{f_{i,j}-f_{\min}}{f_{\max}-f_{\min}}\cdot 255$
		\item Redimensionnement à l'aide de la fonction \texttt{imresize}.
	\end{itemize}
\end{enumerate}

\subsection{Classification multi-classes}

Cette partie de TP est proposée uniquement pour ceux qui auraient déjà tout fait.

Il y a deux stratégies principales pour entraîner un classifieur multi-classe (avec M classes) à partir des outils de classification bi-classes que nous avons vues:

\paragraph{One-vs-one:} on entraine des classifieurs SVM bi-classes pour chacune des  paires de classes. Par exemple : ’1’ vs. ’2’, ’1’ vs. ’3’, ..., ’1’ vs. ’0’, ’2’ vs.’3’, ..., ’2’ vs. ’0’, ... etc. Pour trouver une prediction, on évalue ensuite tous les classifieurs sur un nouvel echantillon, et on choisit la classe qui ressort le plus souvent. Par exemple :
\begin{itemize}[label=-]
	\item '1' vs. '2' : '1'
	\item '1' vs. ’3’ : '1'
	\item '2' vs. '3' : '3'
\end{itemize}
$\Rightarrow$ la classe '1' est choisie.

\paragraph{One-vs-all:} On entraine M classifieurs pour la classification binaire :
\begin{itemize}[label=-]
	\item '1' vs. '2','3',...'0'
	\item '2' vs. '1','3',...'0'
	\item $\vdots$
	\item '0' vs. '1','2',...'9'
\end{itemize}

\begin{enumerate}
	\item Implémenter une des deux stratégies pour le jeu de données USPS. Calculer la matrice de confusion et le taux d'erreur. 
\end{enumerate}